\documentclass[12pt, a4paper]{article}

\usepackage{amsmath, amssymb}

\usepackage{fontspec}
\setmainfont{FreeSerif}

\usepackage{geometry}
\geometry{margin=1in}

\usepackage[colorlinks=true, linkcolor=blue, urlcolor=blue, citecolor=blue]{hyperref}

\setcounter{secnumdepth}{0} 

\begin{document}

\tableofcontents
\vspace{2em}
\hrule
\vspace{2em}

\section{Մատրիցի ձևափոխումը անկյունագծային տեսքի}

\subsection{1. Մատրիցների նմանության տեսություն}

\subsubsection{1.1. Սահմանում}

Երկու $n$-րդ կարգի $A$ և $B$ քառակուսի մատրիցներ կոչվում են \textbf{նման} ($A \sim B$), եթե գոյություն ունի այնպիսի $P$ հակադարձելի մատրից, որ.

$$
B = P A P^{-1}
$$

\begin{quote}
\textbf{Ինտուիտիվ մեկնաբանություն.} Նման մատրիցները ներկայացնում են միևնույն գծային օպերատորը (ձևափոխությունը) տարբեր կոորդինատային համակարգերում (բազիսներում): $P$ մատրիցը հանդես է գալիս որպես «թարգմանիչ» մի բազիսից մյուսը անցնելու համար:
\end{quote}

\subsubsection{1.2. Հակադարձելիության պահանջը}

$P$ անցման մատրիցը պարտադիր պետք է լինի հակադարձելի ($\det(P) \neq 0$), քանի որ.

\begin{itemize}
    \item \textbf{Տեղեկատվության պահպանում.} Այն երաշխավորում է, որ ձևափոխության ընթացքում տեղեկատվության կորուստ տեղի չի ունենում և հնարավոր է իրականացնել հետադարձ թարգմանություն:
    \item \textbf{Գծային անկախություն.} Կոորդինատային առանցքները ներկայացնող վեկտորները պետք է լինեն գծորեն անկախ, ինչը հակադարձելի մատրիցի հիմնարար հատկությունն է:
\end{itemize}

\vspace{1em}
\hrule
\vspace{1em}

\subsection{2. Անկյունագծայնացման գործընթացը}

\subsubsection{2.1. Հիմնական հասկացություն}

Եթե գոյություն ունի այնպիսի $P$ հակադարձելի մատրից, որ

$$
P^{-1} A P = D
$$

մատրիցն անկյունագծային է, ապա ասում են, որ $A$ մատրիցը կարելի է \textbf{բերել անկյունագծային տեսքի}: Այս դեպքում $A$ մատրիցը ներկայացվում է որպես.

$$
A = P D P^{-1}
$$

\subsubsection{2.2. Անկյունագծայնացման առավելությունները}

\begin{itemize}
    \item \textbf{Կապերի քանդում (Decoupling).} Ոչ անկյունագծային մատրիցում փոփոխականները «խճճված» են (coupled): Անկյունագծայնացումը թույլ է տալիս առանձնացնել դրանք այնպես, որ յուրաքանչյուր կոորդինատային առանցքի վրա գործողությունը լինի անկախ:
    \item \textbf{Հաշվողական արդյունավետություն.} $A^n$ հաշվելը բարդ է, սակայն $D^n$ հաշվելը վայրկյանական է, քանի որ անկյունագծի տարրերը պարզապես բարձրացվում են համապատասխան աստիճան:
\end{itemize}

\subsubsection{2.3. Հիմնարար թեորեմներ}

\begin{itemize}
    \item \textbf{Թեորեմ 44.} $A$ ($n \times n$) կարգի մատրիցը կարելի է բերել անկյունագծային տեսքի այն և միայն այն դեպքում, երբ այն ունի $n$ հատ գծորեն անկախ սեփական վեկտորներ:
    \item \textbf{Թեորեմ 45.} Եթե $A$-ն անկյունագծայնացվող է, ապա $P$ մատրիցի սյուները $A$-ի գծորեն անկախ սեփական վեկտորներն են, իսկ $D$ մատրիցի գլխավոր անկյունագծի տարրերը՝ $A$-ի սեփական արժեքները:
    \item \textbf{Թեորեմ 46.} Եթե $A$ ($n \times n$) կարգի մատրիցն ունի $n$ իրարից տարբեր սեփական արժեքներ, ապա այն կարելի է բերել անկյունագծային տեսքի:
\end{itemize}

\vspace{1em}
\hrule
\vspace{1em}

\subsection{3. Տեսական օրինակներ}

\subsubsection{Օրինակ 104.}
$3 \times 3$ մատրիցի անկյունագծայնացում

Տրված է

$$
A =
\begin{pmatrix}
3 & -2 & 0 \\
-2 & 3 & 0 \\
0 & 0 & 5
\end{pmatrix}
$$

\begin{itemize}
    \item \textbf{Սեփական արժեքներ.} $\lambda = 1; 5$:
    \item \textbf{Սեփական վեկտորներ.}
\end{itemize}

$$
p_1 = (1, 1, 0)^T, \quad
p_2 = (-1, 1, 0)^T, \quad
p_3 = (0, 0, 1)^T
$$

\begin{itemize}
    \item \textbf{Վերլուծություն.}
\end{itemize}

$$
P =
\begin{pmatrix}
-1 & 0 & 1 \\
1 & 0 & 1 \\
0 & 1 & 0
\end{pmatrix},
\quad
D =
\begin{pmatrix}
5 & 0 & 0 \\
0 & 5 & 0 \\
0 & 0 & 1
\end{pmatrix}
$$

\subsubsection{Օրինակ 107. Ոչ անկյունագծայնացվող մատրից}

$$
A =
\begin{pmatrix}
-3 & 2 \\
-2 & 1
\end{pmatrix}
$$

մատրիցն ունի $\lambda = -2$ կրկնապատիկ սեփական արժեքը, սակայն սեփական ենթատարածությունը մեկ չափանի է: Քանի որ հնարավոր չէ գտնել 2 գծորեն անկախ վեկտոր, մատրիցը հնարավոր չէ բերել անկյունագծային տեսքի:

\subsubsection{Օրինակ 108.}
$A^n$ հաշվարկը

$$
A =
\begin{pmatrix}
3 & 1 & 0 \\
0 & 1 & 0 \\
4 & 2 & 1
\end{pmatrix}
$$

մատրիցի համար ստացվել է $A = P D P^{-1}$ վերլուծությունը, որտեղից.

$$
A^{100} = P D^{100} P^{-1}=
\begin{pmatrix}
3^{100} & \frac{1}{2}(3^{100} - 1) & 0 \\
0 & 1 & 0 \\
2(3^{100} - 1) & 3^{100} - 1 & 1
\end{pmatrix}
$$

\subsection{4. Գործնական վարժություններ}

\subsection{4.1. Անկյունագծայնացման գործնական քայլերը (Վարժություն 253)}

Մատրիցը անկյունագծային տեսքի բերելու համար անհրաժեշտ է հետևել այս հաջորդականությանը.

\textbf{Բնութագրիչ հավասարման լուծում}

Գտնել

$$
\det(A - \lambda I) = 0
$$

հավասարման արմատները։

\begin{itemize}
    \item \textbf{Օրինակ ա)}
\end{itemize}

$$
A =
\begin{pmatrix}
-14 & 12 \\
-20 & 17
\end{pmatrix}
$$

$$
(-14 - \lambda)(17 - \lambda) - (12 \cdot (-20)) = 0
$$

$$
\lambda^2 - 17\lambda + 14\lambda - 238 + 240 = 0
\Rightarrow
\lambda^2 - 3\lambda + 2 = 0
$$

Արմատներն են՝ $\lambda_1 = 1, \lambda_2 = 2$։

\textbf{Քայլ 2. Սեփական վեկտորների որոշում}

Յուրաքանչյուր $\lambda$-ի համար լուծել

$$
(A - \lambda I)v = 0
$$

համասեռ համակարգը։

\begin{itemize}
    \item \textbf{$\lambda = 1$-ի համար}
\end{itemize}

$$
\begin{pmatrix}
-15 & 12 \\
-20 & 16
\end{pmatrix}
\begin{pmatrix}
x \\
y
\end{pmatrix}
=0
$$

$$
-15x + 12y = 0 \Rightarrow 5x = 4y
$$

Ընտրում ենք նվազագույն ամբողջ թվերը՝

$$
x=4,\; y=5 \Rightarrow v_1 = (4,5)^T
$$

\begin{itemize}
    \item \textbf{$\lambda = 2$-ի համար}
\end{itemize}

$$
\begin{pmatrix}
-16 & 12 \\
-20 & 15
\end{pmatrix}
\begin{pmatrix}
x \\
y
\end{pmatrix}
=0
$$

$$
-16x + 12y = 0 \Rightarrow 4x = 3y
$$

$$
x=3,\; y=4 \Rightarrow v_2 = (3,4)^T
$$

\textbf{Քայլ 3. $P$ և $D$ մատրիցների կառուցում}

$$
P = [v_1 | v_2] =
\begin{pmatrix}
4 & 3 \\
5 & 4
\end{pmatrix}
$$

$$
D = \mathrm{diag}(\lambda_1,\lambda_2) =
\begin{pmatrix}
1 & 0 \\
0 & 2
\end{pmatrix}
$$

\vspace{1em}
\hrule
\vspace{1em}

\subsection{4.2 Ազատ փոփոխականների դեպքը (Վարժություն 253-դ)}

Երբ ունենք պատիկ սեփական արժեքներ, սեփական վեկտորների քանակը որոշվում է ազատ փոփոխականների թվով։

\textbf{Մատրից}

$$
A =
\begin{pmatrix}
2 & 0 & -2 \\
0 & 3 & 0 \\
0 & 0 & 3
\end{pmatrix}
$$

\begin{itemize}
    \item \textbf{Սեփական արժեքներ.} $\lambda_1 = 2, \lambda_2 = 3$ (կրկնակի)։
    \item \textbf{$\lambda = 3$-ի համար}
\end{itemize}

$$
(A - 3I)v = 0
$$

$$
\begin{pmatrix}
-1 & 0 & -2 \\
0 & 0 & 0 \\
0 & 0 & 0
\end{pmatrix}
\begin{pmatrix}
x \\
y \\
z
\end{pmatrix}
=0
$$

$$
-x - 2z = 0 \Rightarrow x = -2z
$$

Քանի որ $y$-ի և $z$-ի վրա սահմանափակում չկա, դրանք \textbf{ազատ փոփոխականներ} են։

\begin{enumerate}
    \item \textbf{Ընտրություն 1.}
\end{enumerate}

$$
y=1,\; z=0 \Rightarrow x=0 \Rightarrow v_2 = (0,1,0)^T
$$

\begin{enumerate}
    \setcounter{enumi}{1}
    \item \textbf{Ընտրություն 2.}
\end{enumerate}

$$
y=0,\; z=1 \Rightarrow x=-2 \Rightarrow v_3 = (-2,0,1)^T
$$

\vspace{1em}
\hrule
\vspace{1em}

\subsection{4.3 Մատրիցի աստիճանի բարձրացում (Վարժություն 256)}

Մեթոդը հիմնված է հետևյալ նույնության վրա.

$$
A^n = P D^n P^{-1}
$$

\textbf{Օրինակ.} Հաշվել $A^{10}$ մատրիցը, եթե.

$$
A = \begin{pmatrix} 1 & 0 \\ -1 & 2 \end{pmatrix}
$$

\vspace{1em}
\hrule
\vspace{1em}

\subsubsection{Քայլ 1. $P^{-1}$ հաշվարկը}

$2 \times 2$ մատրիցի հակադարձը հաշվում ենք $P^{-1} = \frac{1}{\det(P)} \text{adj}(P)$ բանաձևով:

$$
P = \begin{pmatrix} 1 & 0 \\ 1 & 1 \end{pmatrix} \implies \det(P) = 1
$$

$$
P^{-1} = \begin{pmatrix} 1 & 0 \\ -1 & 1 \end{pmatrix}
$$

\subsubsection{Քայլ 2. Աստիճանի բարձրացում}

Անկյունագծային մատրիցի դեպքում բարձրացնում ենք միայն անկյունագծի տարրերը.

$$
D^{10} = \begin{pmatrix} 1^{10} & 0 \\ 0 & 2^{10} \end{pmatrix} = \begin{pmatrix} 1 & 0 \\ 0 & 1024 \end{pmatrix}
$$

\subsubsection{Քայլ 3. Մատրիցային բազմապատկում}

Նախ հաշվում ենք $P \cdot D^{10}$, ապա արդյունքը բազմապատկում $P^{-1}$-ով.

$$
P D^{10} = \begin{pmatrix} 1 & 0 \\ 1 & 1 \end{pmatrix} \begin{pmatrix} 1 & 0 \\ 0 & 1024 \end{pmatrix} = \begin{pmatrix} 1 & 0 \\ 1 & 1024 \end{pmatrix}
$$

\textbf{Վերջնական արդյունք.}

$$
A^{10} = \begin{pmatrix} 1 & 0 \\ 1 & 1024 \end{pmatrix} \begin{pmatrix} 1 & 0 \\ -1 & 1 \end{pmatrix} = \begin{pmatrix} 1 & 0 \\ -1023 & 1024 \end{pmatrix}
$$

\subsection{5. Լրացուցիչ սահմանումներ}

\begin{itemize}
    \item \textbf{Ոչ բացասական մատրից ($A \ge 0$):} Եթե բոլոր $a_{ij} \ge 0$:
    \item \textbf{Դրական մատրից ($A > 0$):} Եթե բոլոր $a_{ij} > 0$:
\end{itemize}

\end{document}